\section{Техническое задание}
\subsection{Основание для разработки}

Основанием для разработки является задание на выпускную квалификационную работу бакалавра «Разработка платформы онлайн-бронирования услуг на базе событийно-ориентированной микросервисной архитектуры».

\subsection{Цель и назначение разработки}

Целью разработки является создание программной платформы для автоматизации процесса записи клиентов на услуги частного мастера через веб-интерфейс с учетом рабочего расписания, занятости и уведомлений.

Назначение разрабатываемой платформы:
\begin{itemize}
\item снижение доли ручных операций при ведении клиентской записи;
\item централизация данных о пользователях, приглашениях и бронированиях;
\item повышение надежности процессов записи, отмены и изменения бронирований;
\item обеспечение интеграции с внешними сервисами календаря и уведомлений.
\end{itemize}

Задачами разработки являются:
\begin{itemize}
\item анализ предметной области онлайн-бронирования услуг и выявление требований к разрабатываемой платформе;
\item проектирование архитектуры программного средства, базы данных и пользовательского интерфейса;
\item разработка серверной и клиентской частей платформы с учетом ролевой модели мастера, клиента и администратора;
\item реализация и интеграция модулей управления расписанием, приглашениями, бронированиями, внешним календарем и уведомлениями;
\item тестирование, подготовка к развертыванию и оценка работоспособности разработанной платформы.
\end{itemize}

\subsection{Требования к данным}

Программная система должна обеспечивать централизованное хранение данных, необходимых для регистрации пользователей, настройки расписания мастера, подключения клиентов, создания бронирований, учета внешней занятости и отправки уведомлений.

В системе должны храниться следующие основные группы данных:
\begin{enumerate}
\item Данные пользователей. Система должна хранить идентификатор пользователя, адрес электронной почты, хеш пароля, основную роль, назначенные роли, сведения Telegram-аккаунта, дату создания и дату последнего изменения учетной записи.
\item Данные профиля мастера. Система должна хранить связь мастера с учетной записью пользователя, отображаемое имя, описание профиля и служебные даты создания и изменения записи.
\item Данные профиля клиента. Система должна хранить связь клиента с учетной записью пользователя, отображаемое имя и служебные даты создания и изменения записи.
\item Данные связей мастера и клиента. Система должна хранить сведения о том, какие клиенты подключены к конкретному мастеру, а также дату установления такой связи.
\item Данные приглашений. Система должна хранить уникальный токен приглашения, мастера, создавшего приглашение, клиента, принявшего приглашение, статус приглашения, дату создания, дату истечения, дату принятия и дату отзыва.
\item Данные расписания и рабочей доступности. Система должна хранить часовой пояс мастера, длительность одного слота, регулярные правила доступности по дням недели, время начала и окончания рабочего интервала, название правила и признак активности.
\item Данные бронирований. Система должна хранить мастера, клиента, статус бронирования, начало и окончание временного интервала, примечание клиента, дату создания, дату изменения и дату отмены записи.
\item Данные внешнего календаря. Система должна хранить сведения о подключении внешнего календаря: провайдера, мастера, идентификатор календаря, email подключенного аккаунта, токены доступа, область доступа, время истечения токена, дату подключения и дату отключения, а также события внешнего календаря и интервалы внешней занятости, связанные с бронированиями и задачами синхронизации.
\item Данные уведомлений. Система должна хранить каналы доставки уведомлений пользователя, тип канала, адрес доставки, отображаемое имя канала, внешний идентификатор, статус активности, а также одноразовые токены привязки Telegram-аккаунта.
\item Служебные данные аудита и событийной публикации. Система должна хранить административные действия, эксплуатационные события и исходящие доменные события, необходимые для последующей публикации в Kafka и доставки уведомлений.
\end{enumerate}

К данным программной системы предъявляются следующие требования:
\begin{itemize}
\item данные должны храниться в реляционной базе данных с соблюдением ссылочной целостности;
\item идентификаторы основных сущностей должны быть уникальными;
\item email пользователя, токен приглашения, идентификатор внешнего события и идентификатор outbox-события не должны дублироваться в пределах соответствующих сущностей;
\item активные бронирования одного мастера не должны пересекаться по времени;
\item временной интервал бронирования должен иметь корректные границы: время начала должно быть меньше времени окончания;
\item отмененные бронирования не должны блокировать повторное использование временного слота;
\item пароли пользователей должны храниться только в виде хеша;
\item токены внешних интеграций не должны отображаться пользователю в интерфейсе и не должны попадать в открытые журналы;
\item данные, используемые для отправки уведомлений, должны позволять повторную обработку события при временной недоступности внешнего канала;
\item изменения критичных данных должны фиксироваться так, чтобы администратор мог восстановить контекст выполненного действия.
\end{itemize}

\subsection{Требования к функциональности}

Разрабатываемая платформа должна обеспечивать выполнение следующих функций:
\begin{enumerate}
\item Управление пользователями. Система должна обеспечивать регистрацию, аутентификацию и авторизацию пользователей с разделением ролей «мастер», «клиент» и «администратор».
\item Управление профилем. Пользователь должен иметь возможность просматривать и редактировать основные данные профиля в рамках своей роли.
\item Управление рабочей доступностью мастера. Мастер должен иметь возможность задавать рабочие интервалы, длительность слота, часовой пояс и актуализировать эти параметры без перезапуска системы.
\item Работа с приглашениями. Мастер должен иметь возможность создавать и контролировать пригласительные ссылки для подключения клиентов.
\item Создание бронирования. Клиент должен иметь возможность выбрать доступный слот и создать запись в пределах действующего приглашения и правил доступа.
\item Изменение и отмена бронирования. Мастер и клиент должны иметь возможность изменять или отменять только принадлежащие им бронирования в соответствии с правами доступа.
\item Расчет доступных слотов. Система должна рассчитывать доступные интервалы с учетом рабочего расписания мастера, уже созданных бронирований и занятости во внешнем календаре.
\item Интеграция с внешним календарем. Платформа должна обеспечивать подключение внешнего календаря, учет занятых интервалов и синхронизацию событий бронирования.
\item Уведомления. Система должна автоматически отправлять уведомления о создании, изменении и отмене бронирований.
\item Журналирование и мониторинг. Система должна фиксировать ключевые действия пользователей и обеспечивать выдачу базовых эксплуатационных метрик для контроля состояния платформы.
\end{enumerate}

\subsection{Требования пользователя к интерфейсу web-приложения}

Интерфейс платформы должен быть реализован в виде веб-приложения и обеспечивать удобную навигацию для мастера, клиента и администратора.

Пользовательский интерфейс должен включать:
\begin{itemize}
\item публичную стартовую страницу с кратким описанием сервиса и переходом к регистрации мастера;
\item формы регистрации и авторизации;
\item личный кабинет мастера;
\item личный кабинет клиента;
\item административный раздел для контроля пользователей, записей, аудита и мониторинга;
\item раздел настройки рабочей доступности мастера;
\item раздел работы с приглашениями;
\item раздел просмотра и управления бронированиями;
\item раздел подключения внешнего календаря;
\item раздел настроек уведомлений, безопасности и профиля.
\end{itemize}

Интерфейс должен быть адаптирован для работы в распространенных браузерах и корректно отображаться на настольных и мобильных устройствах.

Ключевые экраны пользовательского интерфейса представлены на рисунках \ref{landing:image} -- \ref{admin-audit:image}.

На рисунке \ref{landing:image} представлена публичная стартовая страница платформы.

\begin{figure}[H]
\center{\includegraphics[width=0.62\linewidth]{landing}}
\caption{Публичная стартовая страница платформы}
\label{landing:image}
\end{figure}

Как показано на рисунке, стартовая страница кратко объясняет назначение сервиса, разделяет сценарии мастера и клиента, содержит переход к началу работы мастера и блок поддержки. Такой экран нужен до авторизации: он помогает пользователю понять, что мастер публикует слоты и управляет записями, а клиент записывается по ссылке-приглашению.

На рисунке \ref{registration:image} представлена форма регистрации мастера.

\begin{figure}[H]
\center{\includegraphics[width=0.9\linewidth]{registration}}
\caption{Форма регистрации мастера}
\label{registration:image}
\end{figure}

Как показано на рисунке, форма содержит поля ввода email, пароля и подтверждения пароля, кнопку создания учетной записи, переключатель языка интерфейса и ссылку для перехода к входу в существующий аккаунт. После успешной регистрации мастер получает доступ к рабочему пространству для управления слотами и записями.

На рисунке \ref{login:image} представлена форма входа в систему.

\begin{figure}[H]
\center{\includegraphics[width=0.9\linewidth]{login}}
\caption{Форма входа в систему}
\label{login:image}
\end{figure}

Форма входа используется для повторной авторизации уже зарегистрированного пользователя. На экране расположены поля email и пароля, кнопка входа, переключатель языка интерфейса и ссылка для перехода к созданию нового аккаунта.

На рисунке \ref{profile:image} представлен экран профиля пользователя.

\begin{figure}[H]
\center{\includegraphics[width=0.9\linewidth]{profile}}
\caption{Экран профиля пользователя}
\label{profile:image}
\end{figure}

Данный экран содержит сведения о рабочем пространстве мастера. Пользователь может просмотреть email учетной записи, изменить отображаемое имя и описание профиля, сохранить изменения, а также выполнить выход из текущей сессии в отдельном блоке управления сессией.

На рисунке \ref{invite:image} представлен экран работы с клиентами и пригласительной ссылкой.

\begin{figure}[H]
\center{\includegraphics[width=0.9\linewidth]{invite}}
\caption{Экран работы с клиентами и пригласительной ссылкой}
\label{invite:image}
\end{figure}

Экран предназначен для управления постоянной ссылкой-приглашением и просмотра клиентов мастера. В интерфейсе отображается адрес для входа клиента, кнопка копирования ссылки, а также таблица подключенных клиентов с датой присоединения.

На рисунке \ref{slots:image} представлен экран календаря записей мастера.

\begin{figure}[H]
\center{\includegraphics[width=0.9\linewidth]{slots}}
\caption{Экран календаря записей мастера}
\label{slots:image}
\end{figure}

Данный экран позволяет мастеру контролировать рабочее расписание в календарном виде. В интерфейсе отображаются выбранная неделя, часовой пояс, счетчики свободных и занятых интервалов, занятые записи клиентов, боковое меню разделов и элементы переключения между недельным и месячным режимами.

На рисунке \ref{services:image} представлен экран подключения внешних сервисов.

\begin{figure}[H]
\center{\includegraphics[width=0.9\linewidth]{services}}
\caption{Экран подключения внешних сервисов}
\label{services:image}
\end{figure}

Экран подключения внешних сервисов используется для управления календарными провайдерами. В представленном варианте отображается интеграция с Google Calendar, состояние подключения, email подключенного аккаунта и элементы управления для подключения или отключения календаря.

На рисунке \ref{notification:image} представлен экран настроек уведомлений.

\begin{figure}[H]
\center{\includegraphics[width=0.9\linewidth]{notification}}
\caption{Экран настроек уведомлений}
\label{notification:image}
\end{figure}

Экран настроек уведомлений отображает подключенные каналы информирования пользователя. В представленном варианте показан Telegram-канал, его статус, связанный аккаунт и элемент управления для отключения уведомлений.

На рисунке \ref{security:image} представлен экран настроек безопасности.

\begin{figure}[H]
\center{\includegraphics[width=0.9\linewidth]{security}}
\caption{Экран настроек безопасности}
\label{security:image}
\end{figure}

Экран настроек безопасности позволяет пользователю изменить пароль учетной записи. Форма содержит поля текущего и нового пароля, а также отдельную кнопку подтверждения изменения.

На рисунке \ref{client-overview:image} представлена панель клиента.

\begin{figure}[H]
\center{\includegraphics[width=0.9\linewidth]{client-overview}}
\caption{Панель клиента}
\label{client-overview:image}
\end{figure}

Панель клиента используется для просмотра подключенных мастеров и перехода к основным клиентским сценариям. В интерфейсе отображается текущая учетная запись клиента, быстрые переходы к записям, слотам и настройкам аккаунта, а также таблица мастеров, связанных с клиентом через приглашения или созданные записи.

На рисунке \ref{client-appointments:image} представлен экран записей клиента.

\begin{figure}[H]
\center{\includegraphics[width=0.9\linewidth]{client-appointments}}
\caption{Экран записей клиента}
\label{client-appointments:image}
\end{figure}

Экран записей клиента отображает активные и архивные бронирования текущего пользователя. Таблица содержит идентификатор записи, состояние, мастера, время начала и окончания, а фильтры позволяют быстро переключаться между активными и отмененными записями.

На рисунке \ref{client-slots:image} представлен экран бронирования свободного слота клиентом.

\begin{figure}[H]
\center{\includegraphics[width=0.9\linewidth]{client-slots}}
\caption{Экран бронирования свободного слота клиентом}
\label{client-slots:image}
\end{figure}

Экран бронирования позволяет клиенту выбрать подключенного мастера и свободный временной интервал в календаре. В интерфейсе отображаются сведения о мастере, его часовом поясе, выбранной неделе, счетчики свободных и занятых интервалов, а свободный слот содержит действие для создания записи.

На рисунке \ref{admin-overview:image} представлена административная сводка платформы.

\begin{figure}[H]
\center{\includegraphics[width=0.9\linewidth]{admin-overview}}
\caption{Административная сводка платформы}
\label{admin-overview:image}
\end{figure}

Административная сводка отображает текущее состояние платформы: количество активных и отключенных пользователей, число администраторов, активные и отмененные записи, а также график системных итогов. Из этого раздела администратор может перейти к расширенному мониторингу в Grafana.

На рисунке \ref{admin-monitoring:image} представлен экран административного мониторинга.

\begin{figure}[H]
\center{\includegraphics[width=0.9\linewidth]{admin-monitoring}}
\caption{Экран административного мониторинга}
\label{admin-monitoring:image}
\end{figure}

Экран мониторинга показывает состояние Prometheus, количество активных alerts, число недоступных targets, график метрик за выбранный период и таблицу scrape targets. Такой экран нужен для оперативного контроля работоспособности сервисов и инфраструктурных компонентов.

На рисунке \ref{admin-users:image} представлен экран управления пользователями.

\begin{figure}[H]
\center{\includegraphics[width=0.9\linewidth]{admin-users}}
\caption{Экран управления пользователями}
\label{admin-users:image}
\end{figure}

Экран управления пользователями позволяет администратору искать учетные записи, фильтровать их по роли и статусу, просматривать список пользователей и изменять выбранную учетную запись. В карточке пользователя отображаются имя, статус, назначенные роли и связанные профили мастера или клиента.

На рисунке \ref{admin-audit:image} представлен журнал административных действий.

\begin{figure}[H]
\center{\includegraphics[width=0.9\linewidth]{admin-audit}}
\caption{Журнал административных действий}
\label{admin-audit:image}
\end{figure}

Журнал содержит фильтры по пользователю-исполнителю, действию, типу и идентификатору целевого объекта. Таблица показывает время операции, исполнителя, действие, целевой объект и correlation id, что позволяет восстановить контекст ручного вмешательства администратора.

Интерфейс системы должен обеспечивать логичное расположение элементов управления, единообразное оформление экранов, наглядное представление состояния записи и минимальное количество действий пользователя при выполнении основных операций.

\subsection{Требования к обработке ошибок}

Система должна обеспечивать корректную обработку исключительных ситуаций.

Основные требования:
\begin{itemize}
\item при возникновении ошибки пользователю должно отображаться понятное сообщение без технических деталей реализации;
\item внутренние ошибки сервиса не должны раскрывать служебную информацию;
\item ошибки должны фиксироваться в журнале событий с указанием времени, типа ошибки и контекста операции;
\item должны обрабатываться ошибки валидации входных данных;
\item должны обрабатываться ошибки авторизации и доступа к защищенным операциям;
\item должны обрабатываться ошибки интеграции с внешним календарем и модулем уведомлений;
\item при временной недоступности внешнего сервиса операция должна завершаться контролируемым отказом с сохранением целостности данных.
\end{itemize}

\subsection{Требования к удобству использования (UX)}

Система должна обеспечивать удобство и простоту взаимодействия для мастера, клиента и администратора.

Основные требования:
\begin{itemize}
\item интерфейс должен быть интуитивно понятным и содержать логичную структуру разделов;
\item количество действий для выполнения ключевой операции записи должно быть минимальным;
\item состояние бронирований и доступных слотов должно отображаться в наглядном виде;
\item должна быть реализована визуальная обратная связь: индикаторы загрузки, статусы операций, сообщения об успешном выполнении действий;
\item форма ввода данных должна содержать валидацию и подсказки по заполнению;
\item система должна поддерживать корректную работу на настольных и мобильных устройствах.
\end{itemize}

\subsection{Моделирование вариантов использования}

Для проектируемой платформы должна быть сформирована модель вариантов использования на языке UML, обеспечивающая наглядное представление взаимодействия пользователей с системой.

В системе предусмотрены следующие роли пользователей:
\begin{itemize}
\item мастер -- управляет рабочей доступностью, приглашениями и бронированиями;
\item клиент -- присоединяется по приглашению, выбирает слот и управляет своими записями;
\item администратор -- контролирует состояние платформы, пользователей, записи, аудит и эксплуатационные метрики.
\end{itemize}

Основные прецеденты системы:
\begin{enumerate}
\item Регистрация и вход пользователя в систему.
\item Настройка мастером рабочей доступности.
\item Создание мастером пригласительной ссылки.
\item Присоединение клиента по пригласительной ссылке.
\item Просмотр клиентом доступных слотов.
\item Создание клиентом бронирования.
\item Изменение или отмена бронирования участником записи.
\item Получение пользователями уведомлений о событиях записи.
\item Просмотр администратором пользователей, записей, аудита и мониторинга.
\end{enumerate}

На рисунке \ref{client-usecase:image} представлена диаграмма прецедентов клиента.

\begin{figure}[H]
\centering
\resizebox{0.98\linewidth}{!}{%
\begin{tikzpicture}[
usecase/.style={draw, ellipse, align=center, text width=2.9cm, minimum height=0.65cm, font=\scriptsize, inner sep=1pt},
widecase/.style={usecase, text width=3.2cm},
relation/.style={-Latex, line width=0.55pt},
system/.style={draw, rectangle, line width=0.8pt},
block/.style={draw, rectangle, line width=0.55pt}
]
\draw[system] (2.9,0.2) rectangle (14.9,6.7);
\node[font=\bfseries\small] at (8.9,6.45) {Easyappoint};

\draw (1.0,5.5) circle (0.16);
\draw (1.0,5.34) -- (1.0,4.8);
\draw (0.68,5.13) -- (1.32,5.13);
\draw (1.0,4.8) -- (0.7,4.3);
\draw (1.0,4.8) -- (1.3,4.3);
\node[align=center, font=\scriptsize] (guest) at (1.0,3.9) {Неавторизованный\\пользователь};

\draw (1.0,2.45) circle (0.16);
\draw (1.0,2.29) -- (1.0,1.75);
\draw (0.68,2.08) -- (1.32,2.08);
\draw (1.0,1.75) -- (0.7,1.25);
\draw (1.0,1.75) -- (1.3,1.25);
\node[align=center, font=\scriptsize] (client) at (1.0,0.85) {Клиент};

\draw[block] (3.35,4.35) rectangle (7.25,6.05);
\node[font=\bfseries\scriptsize] at (5.3,5.82) {Доступ};
\node[usecase] (reg) at (5.3,5.32) {Регистрация};
\node[usecase] (login) at (5.3,4.65) {Вход в систему};

\draw[block] (3.35,0.55) rectangle (7.75,3.95);
\node[font=\bfseries\scriptsize] at (5.55,3.72) {Запись к мастеру};
\node[widecase] (invite) at (5.55,3.05) {Присоединение\\по приглашению};
\node[widecase] (slots) at (5.55,2.15) {Просмотр\\доступных слотов};
\node[widecase] (booking) at (5.55,1.25) {Создание\\бронирования};

\draw[block] (8.55,0.55) rectangle (14.35,3.95);
\node[font=\bfseries\scriptsize] at (11.45,3.72) {Управление записью};
\node[widecase] (list) at (11.45,3.05) {Просмотр\\своих записей};
\node[widecase] (change) at (11.45,2.15) {Изменение или\\отмена записи};
\node[widecase] (notify) at (11.45,1.25) {Получение\\уведомлений};

\draw[relation] (1.4,5.15) -- (reg.west);
\draw[relation] (1.4,4.95) -- (login.west);
\draw[relation] (1.4,2.15) -- (invite.west);
\draw[relation] (1.4,1.95) -- (slots.west);
\draw[relation] (1.4,1.75) -- (booking.west);
\draw[relation] (booking.east) -- (notify.west);
\draw[relation] (slots.east) -- (list.west);
\end{tikzpicture}%
}
\caption{Диаграмма прецедентов клиента}
\label{client-usecase:image}
\end{figure}

На рисунке \ref{master-usecase:image} представлена диаграмма прецедентов мастера.

\begin{figure}[H]
\centering
\resizebox{0.96\linewidth}{!}{%
\begin{tikzpicture}[
usecase/.style={draw, ellipse, align=center, text width=2.9cm, minimum height=0.65cm, font=\scriptsize, inner sep=1pt},
widecase/.style={usecase, text width=3.25cm},
relation/.style={-Latex, line width=0.55pt},
system/.style={draw, rectangle, line width=0.8pt},
block/.style={draw, rectangle, line width=0.55pt}
]
\draw[system] (2.9,0.2) rectangle (14.4,6.7);
\node[font=\bfseries\small] at (8.65,6.45) {Easyappoint};

\draw (1.0,2.95) circle (0.16);
\draw (1.0,2.79) -- (1.0,2.25);
\draw (0.68,2.58) -- (1.32,2.58);
\draw (1.0,2.25) -- (0.7,1.75);
\draw (1.0,2.25) -- (1.3,1.75);
\node[align=center, font=\scriptsize] (master) at (1.0,1.35) {Мастер};

\draw[block] (3.35,4.35) rectangle (7.25,6.05);
\node[font=\bfseries\scriptsize] at (5.3,5.82) {Доступ};
\node[usecase] (regm) at (5.3,5.32) {Регистрация};
\node[usecase] (loginm) at (5.3,4.65) {Вход в систему};

\draw[block] (3.35,0.55) rectangle (7.9,3.95);
\node[font=\bfseries\scriptsize] at (5.62,3.72) {Рабочее пространство};
\node[widecase] (profile) at (5.62,3.05) {Настройка\\профиля};
\node[widecase] (availability) at (5.62,2.15) {Формирование рабочей\\доступности};
\node[widecase] (invitation) at (5.62,1.25) {Создание\\приглашения};

\draw[block] (8.75,0.55) rectangle (13.95,3.95);
\node[font=\bfseries\scriptsize] at (11.35,3.72) {Записи и уведомления};
\node[widecase] (settings) at (11.35,3.05) {Настройка\\уведомлений};
\node[widecase] (records) at (11.35,2.15) {Просмотр\\бронирований};
\node[widecase] (edit) at (11.35,1.25) {Изменение или\\отмена записи};

\draw[relation] (1.4,2.85) -- (loginm.west);
\draw[relation] (1.4,2.65) -- (profile.west);
\draw[relation] (1.4,2.45) -- (availability.west);
\draw[relation] (1.4,2.25) -- (invitation.west);
\draw[relation] (profile.east) -- (settings.west);
\draw[relation] (availability.east) -- (records.west);
\draw[relation] (invitation.east) -- (edit.west);
\end{tikzpicture}%
}
\caption{Диаграмма прецедентов мастера}
\label{master-usecase:image}
\end{figure}

На рисунке \ref{admin-usecase:image} представлена диаграмма прецедентов администратора.

\begin{figure}[H]
\centering
\resizebox{0.94\linewidth}{!}{%
\begin{tikzpicture}[
usecase/.style={draw, ellipse, align=center, text width=4.0cm, minimum height=0.68cm, font=\scriptsize, inner sep=1pt},
relation/.style={-Latex, line width=0.55pt},
system/.style={draw, rectangle, line width=0.8pt},
block/.style={draw, rectangle, line width=0.55pt}
]
\draw[system] (2.9,0.2) rectangle (14.3,6.2);
\node[font=\bfseries\small] at (8.6,5.92) {Easyappoint};

\draw (1.0,3.25) circle (0.16);
\draw (1.0,3.09) -- (1.0,2.55);
\draw (0.68,2.88) -- (1.32,2.88);
\draw (1.0,2.55) -- (0.7,2.05);
\draw (1.0,2.55) -- (1.3,2.05);
\node[align=center, font=\scriptsize] (admin) at (1.0,1.65) {Администратор};

\draw[block] (4.0,0.75) rectangle (13.35,5.55);
\node[font=\bfseries\scriptsize] at (8.68,5.32) {Административная панель};
\node[usecase] (overview) at (8.68,4.65) {Просмотр сводки};
\node[usecase] (users) at (8.68,3.85) {Просмотр и изменение\\пользователей};
\node[usecase] (appointments) at (8.68,3.05) {Просмотр и изменение\\записей};
\node[usecase] (audit) at (8.68,2.25) {Просмотр аудита};
\node[usecase] (monitoring) at (8.68,1.45) {Просмотр мониторинга};

\draw[relation] (1.4,3.05) -- (overview.west);
\draw[relation] (1.4,2.85) -- (users.west);
\draw[relation] (1.4,2.65) -- (appointments.west);
\draw[relation] (1.4,2.45) -- (audit.west);
\draw[relation] (1.4,2.25) -- (monitoring.west);
\end{tikzpicture}%
}
\caption{Диаграмма прецедентов администратора}
\label{admin-usecase:image}
\end{figure}

\subsubsection{Вариант использования «Вход в систему»}

Заинтересованные лица и их требования: мастер или клиент желает получить доступ к функциям платформы.

Предусловие: пользователь не авторизован.

Постусловие: пользователь авторизован и перенаправлен в соответствующий личный кабинет.

Основной успешный сценарий:
\begin{enumerate}
\item Пользователь открывает страницу входа.
\item Вводит учетные данные.
\item Подтверждает вход.
\item Система выполняет проверку данных.
\item Пользователь получает доступ к системе.
\end{enumerate}

Альтернативный сценарий:
\begin{enumerate}
\item Пользователь вводит некорректные данные.
\item Система отклоняет вход и выводит сообщение об ошибке.
\end{enumerate}

\subsubsection{Вариант использования «Настройка рабочей доступности»}

Заинтересованные лица и их требования: мастер желает указать рабочие интервалы и параметры слотов.

Предусловие: мастер авторизован.

Постусловие: параметры рабочей доступности сохранены и применены к расчету слотов.

Основной успешный сценарий:
\begin{enumerate}
\item Мастер открывает раздел доступности.
\item Задает рабочие интервалы, длительность слота и часовой пояс.
\item Подтверждает сохранение.
\item Система сохраняет настройки и обновляет расчет доступных интервалов.
\end{enumerate}

Альтернативный сценарий:
\begin{enumerate}
\item Введены некорректные параметры.
\item Система отменяет сохранение и выводит сообщение об ошибке.
\end{enumerate}

\subsubsection{Вариант использования «Создание пригласительной ссылки»}

Заинтересованные лица и их требования: мастер желает предоставить клиенту доступ к записи.

Предусловие: мастер авторизован.

Постусловие: сформирована активная пригласительная ссылка.

Основной успешный сценарий:
\begin{enumerate}
\item Мастер открывает раздел приглашений.
\item Нажимает кнопку создания приглашения.
\item Система формирует ссылку и отображает её для копирования.
\end{enumerate}

Альтернативный сценарий:
\begin{enumerate}
\item Возникает ошибка сохранения приглашения.
\item Система выводит сообщение о невозможности выполнения операции.
\end{enumerate}

\subsubsection{Вариант использования «Присоединение клиента по приглашению»}

Заинтересованные лица и их требования: клиент желает подключиться к мастеру по ссылке.

Предусловие: у клиента есть действующая пригласительная ссылка.

Постусловие: клиент связан с мастером и получает доступ к просмотру слотов.

Основной успешный сценарий:
\begin{enumerate}
\item Клиент переходит по ссылке приглашения.
\item Система проверяет корректность и актуальность приглашения.
\item Клиент подтверждает присоединение.
\item Система фиксирует связь клиента с мастером.
\end{enumerate}

Альтернативный сценарий:
\begin{enumerate}
\item Приглашение недействительно или истекло.
\item Система выводит сообщение об ошибке доступа.
\end{enumerate}

\subsubsection{Вариант использования «Просмотр доступных слотов»}

Заинтересованные лица и их требования: клиент желает увидеть свободные интервалы мастера.

Предусловие: клиент присоединен к мастеру.

Постусловие: клиент получает список доступных слотов на выбранный период.

Основной успешный сценарий:
\begin{enumerate}
\item Клиент открывает раздел доступных слотов.
\item Указывает период просмотра.
\item Система рассчитывает и отображает свободные интервалы.
\end{enumerate}

Альтернативный сценарий:
\begin{enumerate}
\item На выбранный период отсутствуют свободные интервалы.
\item Система отображает сообщение об отсутствии доступных слотов.
\end{enumerate}

\subsubsection{Вариант использования «Создание бронирования»}

Заинтересованные лица и их требования: клиент желает забронировать выбранный слот.

Предусловие: выбран доступный слот.

Постусловие: бронирование создано и сохранено в системе.

Основной успешный сценарий:
\begin{enumerate}
\item Клиент выбирает свободный слот.
\item Подтверждает создание записи.
\item Система выполняет проверку доступности слота.
\item Система сохраняет бронирование и формирует событие для уведомлений.
\end{enumerate}

Альтернативный сценарий:
\begin{enumerate}
\item На момент подтверждения слот уже недоступен.
\item Система отклоняет создание бронирования и предлагает выбрать другой интервал.
\end{enumerate}

\subsubsection{Вариант использования «Изменение или отмена бронирования»}

Заинтересованные лица и их требования: мастер или клиент желает изменить или отменить запись.

Предусловие: пользователь имеет права на выбранное бронирование.

Постусловие: состояние бронирования обновлено.

Основной успешный сценарий:
\begin{enumerate}
\item Пользователь открывает список своих записей.
\item Выбирает запись для изменения или отмены.
\item Подтверждает действие.
\item Система обновляет запись и формирует событие для уведомлений.
\end{enumerate}

Альтернативный сценарий:
\begin{enumerate}
\item Пользователь пытается изменить чужое бронирование.
\item Система запрещает выполнение операции.
\end{enumerate}

\subsubsection{Вариант использования «Получение уведомлений о событиях записи»}

Заинтересованные лица и их требования: мастер и клиент желают своевременно получать уведомления о ключевых изменениях записи.

Предусловие: в системе произошло событие создания, изменения или отмены бронирования.

Постусловие: пользователям доставлено уведомление по настроенному каналу.

Основной успешный сценарий:
\begin{enumerate}
\item Система фиксирует событие записи.
\item Событие передается в модуль уведомлений.
\item Формируется уведомление для участников записи.
\item Уведомление доставляется получателям.
\end{enumerate}

Альтернативный сценарий:
\begin{enumerate}
\item Внешний канал уведомлений временно недоступен.
\item Система фиксирует ошибку и выполняет повторную попытку доставки.
\end{enumerate}

\subsubsection{Вариант использования «Административный контроль платформы»}

Заинтересованные лица и их требования: администратор желает контролировать состояние платформы, пользователей, записи, аудит и эксплуатационные метрики.

Предусловие: пользователь авторизован с ролью администратора.

Постусловие: администратор получил актуальные сведения о состоянии платформы; при изменении пользователя или записи система сохранила результат и зафиксировала событие аудита.

Основной успешный сценарий:
\begin{enumerate}
\item Администратор открывает административный раздел.
\item Выбирает сводку, пользователей, записи, аудит или мониторинг.
\item Система проверяет наличие административной роли.
\item Система отображает данные с учетом выбранных фильтров.
\item При необходимости администратор изменяет пользователя или запись.
\item Система сохраняет изменения и фиксирует административное событие в журнале аудита.
\end{enumerate}

Альтернативный сценарий:
\begin{enumerate}
\item Пользователь без роли администратора пытается открыть административный раздел.
\item Система запрещает доступ к административным функциям.
\end{enumerate}

\subsection{Требования к надёжности}

Платформа должна обеспечивать стабильную и корректную работу при штатной эксплуатации.

Основные требования:
\begin{itemize}
\item операции создания, изменения и отмены бронирования должны выполняться атомарно;
\item при сбое внешних интеграций данные в основной базе должны сохранять целостность;
\item система должна корректно восстанавливаться после перезапуска сервисов;
\item критические ошибки и сбои должны фиксироваться в журнале событий;
\item должны быть предусмотрены резервное копирование и восстановление данных.
\end{itemize}

\subsection{Требования к информационной безопасности}

Система должна обеспечивать защиту данных и разграничение прав доступа пользователей.

Основные требования:
\begin{itemize}
\item доступ к защищенным операциям возможен только после успешной аутентификации;
\item хранение учетных данных должно выполняться в защищенном виде (хеширование паролей);
\item проверка ролей и прав доступа должна выполняться на серверной стороне;
\item при обмене данными должны использоваться защищенные сетевые протоколы;
\item необходимо вести журнал событий безопасности (вход, неуспешные попытки доступа, критические операции).
\end{itemize}

\subsection{Требования к программной и технической совместимости}

Система должна быть совместима с типовой современной программно-аппаратной средой.

Требования к серверной среде:
\begin{itemize}
\item операционная система Linux;
\item контейнеризация сервисов с использованием Docker;
\item возможность оркестрации в Kubernetes для целевой среды развертывания;
\item в качестве СУБД~--- PostgreSQL;
\item брокер сообщений Kafka.
\end{itemize}

Требования к клиентской среде:
\begin{itemize}
\item современные веб-браузеры: Google Chrome, Mozilla Firefox, Microsoft Edge, Яндекс Браузер;
\item доступ к сети Интернет для работы с веб-приложением.
\end{itemize}

\subsection{Требования к оформлению документации}

Разработка программной документации и программного изделия должна выполняться в соответствии с требованиями действующих нормативных документов ГОСТ 19.102-77 и ГОСТ 34.601-90. Единая система программной документации.
